\documentclass[tikz,border=4pt]{standalone}
\usepackage{tikz}
\usepackage{polymer-paper-preamble}
\usetikzlibrary{arrows.meta,decorations.pathmorphing}
%% polymer_chain.snippet.tex
%% Inverse vulcanization sulfur-rich linear polymer — ball-and-stick TikZ schematic.
%% Fixed pattern: C₄-S₄-C₄-S₆-C₄-S₄-C₃ (chemistry of sulfur polymer trap sites).
%% License: snippet code MIT-style.
%% Author origin: figure-agent v0.3 — L3 snippet (ball-and-stick unified, reference-grounded).
%%
%% Usage:
%%   %% polymer_chain.snippet.tex
%% Inverse vulcanization sulfur-rich linear polymer — ball-and-stick TikZ schematic.
%% Fixed pattern: C₄-S₄-C₄-S₆-C₄-S₄-C₃ (chemistry of sulfur polymer trap sites).
%% License: snippet code MIT-style.
%% Author origin: figure-agent v0.3 — L3 snippet (ball-and-stick unified, reference-grounded).
%%
%% Usage:
%%   %% polymer_chain.snippet.tex
%% Inverse vulcanization sulfur-rich linear polymer — ball-and-stick TikZ schematic.
%% Fixed pattern: C₄-S₄-C₄-S₆-C₄-S₄-C₃ (chemistry of sulfur polymer trap sites).
%% License: snippet code MIT-style.
%% Author origin: figure-agent v0.3 — L3 snippet (ball-and-stick unified, reference-grounded).
%%
%% Usage:
%%   \input{styles/snippets/polymer_chain.snippet.tex}
%%   \PolymerChainBallStick{x_anchor}{y_anchor}
%%
%% Args:
%%   #1 = anchor_x (TikZ cm, left end position)
%%   #2 = anchor_y (TikZ cm, chain center y)
%%
%% Visuals (Pyun-style ball-and-stick):
%%   - C atom (organic spacer): black filled circle (0.08cm radius, outline 0.28pt)
%%   - S atom: cAmber filled circle (0.10cm radius, outline 0.28pt)
%%   - C-C bond: black line (0.70pt)
%%   - S-S bond: cAmber line (1.0pt)
%%   - bond length: 0.40cm
%%   - Horizontal span ≈ 5.4cm (for 3-chain stacking with moderate gap)

\makeatletter
\newcommand{\PolymerChainBallStick}[2]{%
  \begingroup
  \pgfmathsetmacro{\xstart}{#1}%
  \pgfmathsetmacro{\yc}{#2}%
  \pgfmathsetmacro{\bondlen}{0.40}%
  \pgfmathsetmacro{\radiusC}{0.080}%
  \pgfmathsetmacro{\radiusS}{0.100}%

  % Helper: draw C atom (black ball) at given position
  \newcommand{\DrawC}[1]{%
    \draw[fill=black, line width=0.28pt, draw=black] (#1,\yc) circle (\radiusC);%
  }%
  % Helper: draw S atom (cAmber ball) at given position
  \newcommand{\DrawS}[1]{%
    \draw[fill=cAmber, line width=0.28pt, draw=cAmber!50] (#1,\yc) circle (\radiusS);%
  }%
  % Helper: draw C-C bond (black line)
  \newcommand{\BondCC}[2]{%
    \draw[line width=0.70pt, black] (#1,\yc) -- (#2,\yc);%
  }%
  % Helper: draw S-S bond (cAmber line)
  \newcommand{\BondSS}[2]{%
    \draw[line width=1.0pt, cAmber] (#1,\yc) -- (#2,\yc);%
  }%

  % Organic spacer C₁-C₂-C₃-C₄ (x = 0, 0.4, 0.8, 1.2)
  \DrawC{0.00}%
  \BondCC{0.00}{0.40}%
  \DrawC{0.40}%
  \BondCC{0.40}{0.80}%
  \DrawC{0.80}%
  \BondCC{0.80}{1.20}%
  \DrawC{1.20}%

  % S₁-S₂-S₃-S₄ run (x = 1.6, 2.0, 2.4, 2.8)
  \BondSS{1.20}{1.60}%
  \DrawS{1.60}%
  \BondSS{1.60}{2.00}%
  \DrawS{2.00}%
  \BondSS{2.00}{2.40}%
  \DrawS{2.40}%
  \BondSS{2.40}{2.80}%
  \DrawS{2.80}%

  % Organic spacer C₅-C₆-C₇-C₈ (x = 3.2, 3.6, 4.0, 4.4)
  \BondCC{2.80}{3.20}%
  \DrawC{3.20}%
  \BondCC{3.20}{3.60}%
  \DrawC{3.60}%
  \BondCC{3.60}{4.00}%
  \DrawC{4.00}%
  \BondCC{4.00}{4.40}%
  \DrawC{4.40}%

  % S₅-S₆-S₇-S₈-S₉-S₁₀ run (x = 4.8, 5.2, 5.6, 6.0, 6.4, 6.8) — TRAP SITE
  \BondSS{4.40}{4.80}%
  \DrawS{4.80}%
  \BondSS{4.80}{5.20}%
  \DrawS{5.20}%
  \BondSS{5.20}{5.60}%
  \DrawS{5.60}%
  \BondSS{5.60}{6.00}%
  \DrawS{6.00}%
  \BondSS{6.00}{6.40}%
  \DrawS{6.40}%
  \BondSS{6.40}{6.80}%
  \DrawS{6.80}%

  % Organic spacer C₉-C₁₀-C₁₁-C₁₂ (x = 7.2, 7.6, 8.0, 8.4)
  \BondCC{6.80}{7.20}%
  \DrawC{7.20}%
  \BondCC{7.20}{7.60}%
  \DrawC{7.60}%
  \BondCC{7.60}{8.00}%
  \DrawC{8.00}%
  \BondCC{8.00}{8.40}%
  \DrawC{8.40}%

  % S₁₁-S₁₂-S₁₃-S₁₄ run (x = 8.8, 9.2, 9.6, 10.0)
  \BondSS{8.40}{8.80}%
  \DrawS{8.80}%
  \BondSS{8.80}{9.20}%
  \DrawS{9.20}%
  \BondSS{9.20}{9.60}%
  \DrawS{9.60}%
  \BondSS{9.60}{10.00}%
  \DrawS{10.00}%

  % Organic terminal C₁₃-C₁₄-C₁₅ (x = 10.4, 10.8, 11.2)
  \BondCC{10.00}{10.40}%
  \DrawC{10.40}%
  \BondCC{10.40}{10.80}%
  \DrawC{10.80}%
  \BondCC{10.80}{11.20}%
  \DrawC{11.20}%

  \endgroup
}
\makeatother

%%   \PolymerChainBallStick{x_anchor}{y_anchor}
%%
%% Args:
%%   #1 = anchor_x (TikZ cm, left end position)
%%   #2 = anchor_y (TikZ cm, chain center y)
%%
%% Visuals (Pyun-style ball-and-stick):
%%   - C atom (organic spacer): black filled circle (0.08cm radius, outline 0.28pt)
%%   - S atom: cAmber filled circle (0.10cm radius, outline 0.28pt)
%%   - C-C bond: black line (0.70pt)
%%   - S-S bond: cAmber line (1.0pt)
%%   - bond length: 0.40cm
%%   - Horizontal span ≈ 5.4cm (for 3-chain stacking with moderate gap)

\makeatletter
\newcommand{\PolymerChainBallStick}[2]{%
  \begingroup
  \pgfmathsetmacro{\xstart}{#1}%
  \pgfmathsetmacro{\yc}{#2}%
  \pgfmathsetmacro{\bondlen}{0.40}%
  \pgfmathsetmacro{\radiusC}{0.080}%
  \pgfmathsetmacro{\radiusS}{0.100}%

  % Helper: draw C atom (black ball) at given position
  \newcommand{\DrawC}[1]{%
    \draw[fill=black, line width=0.28pt, draw=black] (#1,\yc) circle (\radiusC);%
  }%
  % Helper: draw S atom (cAmber ball) at given position
  \newcommand{\DrawS}[1]{%
    \draw[fill=cAmber, line width=0.28pt, draw=cAmber!50] (#1,\yc) circle (\radiusS);%
  }%
  % Helper: draw C-C bond (black line)
  \newcommand{\BondCC}[2]{%
    \draw[line width=0.70pt, black] (#1,\yc) -- (#2,\yc);%
  }%
  % Helper: draw S-S bond (cAmber line)
  \newcommand{\BondSS}[2]{%
    \draw[line width=1.0pt, cAmber] (#1,\yc) -- (#2,\yc);%
  }%

  % Organic spacer C₁-C₂-C₃-C₄ (x = 0, 0.4, 0.8, 1.2)
  \DrawC{0.00}%
  \BondCC{0.00}{0.40}%
  \DrawC{0.40}%
  \BondCC{0.40}{0.80}%
  \DrawC{0.80}%
  \BondCC{0.80}{1.20}%
  \DrawC{1.20}%

  % S₁-S₂-S₃-S₄ run (x = 1.6, 2.0, 2.4, 2.8)
  \BondSS{1.20}{1.60}%
  \DrawS{1.60}%
  \BondSS{1.60}{2.00}%
  \DrawS{2.00}%
  \BondSS{2.00}{2.40}%
  \DrawS{2.40}%
  \BondSS{2.40}{2.80}%
  \DrawS{2.80}%

  % Organic spacer C₅-C₆-C₇-C₈ (x = 3.2, 3.6, 4.0, 4.4)
  \BondCC{2.80}{3.20}%
  \DrawC{3.20}%
  \BondCC{3.20}{3.60}%
  \DrawC{3.60}%
  \BondCC{3.60}{4.00}%
  \DrawC{4.00}%
  \BondCC{4.00}{4.40}%
  \DrawC{4.40}%

  % S₅-S₆-S₇-S₈-S₉-S₁₀ run (x = 4.8, 5.2, 5.6, 6.0, 6.4, 6.8) — TRAP SITE
  \BondSS{4.40}{4.80}%
  \DrawS{4.80}%
  \BondSS{4.80}{5.20}%
  \DrawS{5.20}%
  \BondSS{5.20}{5.60}%
  \DrawS{5.60}%
  \BondSS{5.60}{6.00}%
  \DrawS{6.00}%
  \BondSS{6.00}{6.40}%
  \DrawS{6.40}%
  \BondSS{6.40}{6.80}%
  \DrawS{6.80}%

  % Organic spacer C₉-C₁₀-C₁₁-C₁₂ (x = 7.2, 7.6, 8.0, 8.4)
  \BondCC{6.80}{7.20}%
  \DrawC{7.20}%
  \BondCC{7.20}{7.60}%
  \DrawC{7.60}%
  \BondCC{7.60}{8.00}%
  \DrawC{8.00}%
  \BondCC{8.00}{8.40}%
  \DrawC{8.40}%

  % S₁₁-S₁₂-S₁₃-S₁₄ run (x = 8.8, 9.2, 9.6, 10.0)
  \BondSS{8.40}{8.80}%
  \DrawS{8.80}%
  \BondSS{8.80}{9.20}%
  \DrawS{9.20}%
  \BondSS{9.20}{9.60}%
  \DrawS{9.60}%
  \BondSS{9.60}{10.00}%
  \DrawS{10.00}%

  % Organic terminal C₁₃-C₁₄-C₁₅ (x = 10.4, 10.8, 11.2)
  \BondCC{10.00}{10.40}%
  \DrawC{10.40}%
  \BondCC{10.40}{10.80}%
  \DrawC{10.80}%
  \BondCC{10.80}{11.20}%
  \DrawC{11.20}%

  \endgroup
}
\makeatother

%%   \PolymerChainBallStick{x_anchor}{y_anchor}
%%
%% Args:
%%   #1 = anchor_x (TikZ cm, left end position)
%%   #2 = anchor_y (TikZ cm, chain center y)
%%
%% Visuals (Pyun-style ball-and-stick):
%%   - C atom (organic spacer): black filled circle (0.08cm radius, outline 0.28pt)
%%   - S atom: cAmber filled circle (0.10cm radius, outline 0.28pt)
%%   - C-C bond: black line (0.70pt)
%%   - S-S bond: cAmber line (1.0pt)
%%   - bond length: 0.40cm
%%   - Horizontal span ≈ 5.4cm (for 3-chain stacking with moderate gap)

\makeatletter
\newcommand{\PolymerChainBallStick}[2]{%
  \begingroup
  \pgfmathsetmacro{\xstart}{#1}%
  \pgfmathsetmacro{\yc}{#2}%
  \pgfmathsetmacro{\bondlen}{0.40}%
  \pgfmathsetmacro{\radiusC}{0.080}%
  \pgfmathsetmacro{\radiusS}{0.100}%

  % Helper: draw C atom (black ball) at given position
  \newcommand{\DrawC}[1]{%
    \draw[fill=black, line width=0.28pt, draw=black] (#1,\yc) circle (\radiusC);%
  }%
  % Helper: draw S atom (cAmber ball) at given position
  \newcommand{\DrawS}[1]{%
    \draw[fill=cAmber, line width=0.28pt, draw=cAmber!50] (#1,\yc) circle (\radiusS);%
  }%
  % Helper: draw C-C bond (black line)
  \newcommand{\BondCC}[2]{%
    \draw[line width=0.70pt, black] (#1,\yc) -- (#2,\yc);%
  }%
  % Helper: draw S-S bond (cAmber line)
  \newcommand{\BondSS}[2]{%
    \draw[line width=1.0pt, cAmber] (#1,\yc) -- (#2,\yc);%
  }%

  % Organic spacer C₁-C₂-C₃-C₄ (x = 0, 0.4, 0.8, 1.2)
  \DrawC{0.00}%
  \BondCC{0.00}{0.40}%
  \DrawC{0.40}%
  \BondCC{0.40}{0.80}%
  \DrawC{0.80}%
  \BondCC{0.80}{1.20}%
  \DrawC{1.20}%

  % S₁-S₂-S₃-S₄ run (x = 1.6, 2.0, 2.4, 2.8)
  \BondSS{1.20}{1.60}%
  \DrawS{1.60}%
  \BondSS{1.60}{2.00}%
  \DrawS{2.00}%
  \BondSS{2.00}{2.40}%
  \DrawS{2.40}%
  \BondSS{2.40}{2.80}%
  \DrawS{2.80}%

  % Organic spacer C₅-C₆-C₇-C₈ (x = 3.2, 3.6, 4.0, 4.4)
  \BondCC{2.80}{3.20}%
  \DrawC{3.20}%
  \BondCC{3.20}{3.60}%
  \DrawC{3.60}%
  \BondCC{3.60}{4.00}%
  \DrawC{4.00}%
  \BondCC{4.00}{4.40}%
  \DrawC{4.40}%

  % S₅-S₆-S₇-S₈-S₉-S₁₀ run (x = 4.8, 5.2, 5.6, 6.0, 6.4, 6.8) — TRAP SITE
  \BondSS{4.40}{4.80}%
  \DrawS{4.80}%
  \BondSS{4.80}{5.20}%
  \DrawS{5.20}%
  \BondSS{5.20}{5.60}%
  \DrawS{5.60}%
  \BondSS{5.60}{6.00}%
  \DrawS{6.00}%
  \BondSS{6.00}{6.40}%
  \DrawS{6.40}%
  \BondSS{6.40}{6.80}%
  \DrawS{6.80}%

  % Organic spacer C₉-C₁₀-C₁₁-C₁₂ (x = 7.2, 7.6, 8.0, 8.4)
  \BondCC{6.80}{7.20}%
  \DrawC{7.20}%
  \BondCC{7.20}{7.60}%
  \DrawC{7.60}%
  \BondCC{7.60}{8.00}%
  \DrawC{8.00}%
  \BondCC{8.00}{8.40}%
  \DrawC{8.40}%

  % S₁₁-S₁₂-S₁₃-S₁₄ run (x = 8.8, 9.2, 9.6, 10.0)
  \BondSS{8.40}{8.80}%
  \DrawS{8.80}%
  \BondSS{8.80}{9.20}%
  \DrawS{9.20}%
  \BondSS{9.20}{9.60}%
  \DrawS{9.60}%
  \BondSS{9.60}{10.00}%
  \DrawS{10.00}%

  % Organic terminal C₁₃-C₁₄-C₁₅ (x = 10.4, 10.8, 11.2)
  \BondCC{10.00}{10.40}%
  \DrawC{10.40}%
  \BondCC{10.40}{10.80}%
  \DrawC{10.80}%
  \BondCC{10.80}{11.20}%
  \DrawC{11.20}%

  \endgroup
}
\makeatother

% Reusable Panel F apparatus motif. The caller owns labels, force arrows,
% panel framing, and whole-panel composition.
\newcommand{\PanelFFloatingCantilever}[1]{%
  \begin{scope}[shift={(#1)}]
    % Stable anchors for semantic overlays and downstream inspection.
    \coordinate (panelFFixedBoundary) at (0,0.34);
    \coordinate (panelFFloatingCantilever) at (-0.26,0);
    \coordinate (panelFDrivenElectrode) at (1.34,-1.40);
    \coordinate (panelFSourceReturn) at (1.74,0.39);
    \coordinate (panelFTrappedCharge) at (-0.22,-0.56);

    % Fixed mechanical boundary and shallow root jaw.
    \draw[cGray!62!black, line width=0.72pt, line cap=round, line join=round]
      (-0.26,0.34) -- (0.26,0.34);
    \foreach \hx in {-0.22,-0.08,0.06,0.20} {
      \draw[cGray!48!black, line width=0.30pt]
        (\hx,0.34) -- ({\hx+0.07},0.42);
    }
    \draw[cGray!62!black, line width=0.72pt, line cap=round]
      (0,0.34) -- (0,0.14);
    \fill[cGray!12, rounded corners=0.22mm]
      (-0.16,0) rectangle (0.16,0.14);
    \draw[cGray!62!black, line width=0.72pt, rounded corners=0.22mm]
      (-0.16,0) rectangle (0.16,0.14);
    \draw[cGray!56!black, line width=0.34pt] (-0.12,0) -- (-0.12,-0.06);
    \draw[cGray!56!black, line width=0.34pt] (0.12,0) -- (0.12,-0.06);

    % Electrically floating charge-trapping cantilever.
    \shade[top color=cAmber!18, bottom color=cAmber!46, rounded corners=0.5mm]
      (-0.14,0) .. controls (-0.02,-0.74) and (-0.49,-1.54) ..
      (-1.37,-2.04) -- (-1.22,-2.20) .. controls (-0.28,-1.72) and
      (0.24,-0.78) .. (0.14,0) -- cycle;
    \draw[cAmber!82!black, line width=0.62pt, rounded corners=0.5mm]
      (-0.14,0) .. controls (-0.02,-0.74) and (-0.49,-1.54) ..
      (-1.37,-2.04) -- (-1.22,-2.20) .. controls (-0.28,-1.72) and
      (0.24,-0.78) .. (0.14,0) -- cycle;
    \draw[cAmber!60!black, line width=0.28pt, opacity=0.42]
      (-0.06,-0.16) .. controls (-0.02,-0.76) and (-0.56,-1.54) .. (-1.29,-2.06);
    \draw[cAmber!60!black, line width=0.28pt, opacity=0.35]
      (0.04,-0.12) .. controls (0.14,-0.82) and (-0.39,-1.72) .. (-1.16,-2.16);

    % Driven electrode and gap guides.
    \shade[left color=cGray!30, right color=cGray!16]
      (1.34,-2.38) rectangle (1.58,-0.02);
    \draw[cGray!82!black, line width=0.62pt]
      (1.34,-2.38) rectangle (1.58,-0.02);
    \draw[cGray!45!black, line width=0.25pt] (1.39,-2.34) -- (1.39,-0.06);
    \foreach \hy in {-2.22,-1.98,-1.74,-1.50,-1.26,-1.02,-0.78,-0.54,-0.30} {
      \draw[cGray!48!black, line width=0.20pt]
        (1.58,\hy) -- (1.34,{\hy-0.05});
    }
    \foreach \yy/\xstart in {-1.88/-0.96,-1.50/-0.67,-1.12/-0.44,-0.74/-0.22} {
      \draw[cGray!23, line width=0.30pt, dash pattern=on 1.3pt off 1.2pt]
        (\xstart,\yy) -- (1.34,\yy);
    }
    \draw[cGray!19, line width=0.22pt, dash pattern=on 1.2pt off 1.5pt]
      (1.08,-0.42) -- (1.34,-0.42);

    % Driven source lead and grounded source return. No lead touches sample.
    \fill[cBlue!5, rounded corners=0.55mm] (0.76,0.40) rectangle (1.24,0.78);
    \draw[cBlue!64!black, line width=0.38pt, rounded corners=0.55mm]
      (0.76,0.40) rectangle (1.24,0.78);
    \node[font=\sffamily\bfseries\fontsize{4.4}{5.2}\selectfont,
          text=cBlue!70!black] at (1.00,0.59) {HV};
    \fill[cBlue!70!black] (1.24,0.50) circle (0.025);
    \draw[cBlue!70!black, line width=0.36pt, rounded corners=0.3mm]
      (1.24,0.50) -- (1.46,0.50) -- (1.46,-0.02);
    \fill[cBlue!70!black] (1.24,0.68) circle (0.025);
    \draw[cBlue!58!black, line width=0.34pt, rounded corners=0.3mm]
      (1.24,0.68) -- (1.74,0.68) -- (1.74,0.39);
    \draw[cBlue!58!black, line width=0.34pt] (1.61,0.39) -- (1.87,0.39);
    \draw[cBlue!58!black, line width=0.34pt] (1.65,0.33) -- (1.83,0.33);
    \draw[cBlue!58!black, line width=0.34pt] (1.69,0.27) -- (1.79,0.27);

    % Flat trapped-charge markers.
    \foreach \cx/\cy/\rr in {-0.22/-0.56/0.075,-0.41/-0.98/0.082,-0.74/-1.48/0.088,-1.12/-1.90/0.082} {
      \fill[cRed!68!black] (\cx,\cy) circle (\rr);
      \draw[cRed!88!black, line width=0.25pt] (\cx,\cy) circle (\rr);
    }
  \end{scope}%
}


\begin{document}
\begin{tikzpicture}[
  font=\sffamily\fontsize{6.2}{7.5}\selectfont,
  panel/.style={font=\bfseries\fontsize{8}{9.5}\selectfont, text=cGray!90!black},
  title/.style={font=\bfseries\fontsize{6.4}{7.7}\selectfont, text=cGray!90!black},
  label/.style={font=\fontsize{5.3}{6.4}\selectfont, text=cGray!82!black, align=center},
  micro/.style={font=\fontsize{4.7}{5.7}\selectfont, text=cGray!70!black, align=center},
  outline/.style={draw=cGray!82!black, line width=0.55pt, line join=round, line cap=round},
  guide/.style={draw=cLGray!80!white, line width=0.35pt},
  axis/.style={-{Stealth[length=1.2mm,width=0.85mm]}, cGray!75!black, line width=0.48pt},
  signal/.style={-{Stealth[length=1.1mm,width=0.8mm]}, cBlue!78!black, line width=0.70pt},
  force/.style={-{Stealth[length=1.3mm,width=0.95mm]}, cRed!78!black, line width=0.85pt},
  bond/.style={cAmber!88!black, line width=0.9pt, line cap=round},
]

% ============================== A: material identity ==============================
\node[panel, anchor=west] at (0.12,11.15) {a};
\node[title, anchor=west] at (0.52,11.15) {sulfur-rich polymer};
\node[label, anchor=west, text=cAmber!75!black] at (0.52,10.68) {schematic sulfur-rich backbone motif};
\PolymerChain{0.75}{9.48}{11}{2,4,6,8,10}
\draw[outline, fill=cAmber!5, rounded corners=1.2mm] (0.72,6.98) rectangle (5.15,8.45);
\PolymerChain{0.92}{8.08}{10}{2,5,8}
\PolymerChain{0.92}{7.67}{10}{3,4,5,7,9}
\PolymerChain{0.92}{7.26}{10}{2,6,8}
\node[label] at (2.90,6.58) {crosslinked sulfur-rich\\polymer network};
\draw[cAmber!65!black, line width=0.42pt] (2.90,6.91) -- (2.90,6.98);

% ============================== B: composition ==============================
\node[panel, anchor=west] at (5.98,11.15) {b};
\node[title, anchor=west] at (6.38,11.15) {composition tunes trap landscape};
\draw[axis] (6.45,6.85) -- (11.20,6.85);
\node[micro, anchor=north] at (11.20,6.75) {sulfur content};
\foreach \x/\tag/\col in {6.95/S60/cBlue,8.78/S75/cTeal,10.60/S85/cRed}{
  \node[font=\bfseries\fontsize{5.7}{6.8}\selectfont, text=\col!80!black] at (\x,10.08) {\tag};
  \draw[\col!78!black, line width=0.85pt] plot[smooth] coordinates {(\x-0.62,8.92) (\x-0.32,9.28) (\x,8.95) (\x+0.32,9.27) (\x+0.62,8.94)};
  \foreach \dx/\dy in {-0.32/0.25,0.00/-0.12,0.32/0.24}{
    \node[circle, fill=\col!70, draw=\col!85!black, minimum size=1.35mm, inner sep=0pt, line width=0.25pt] at (\x+\dx,8.96+\dy) {};
  }
  \draw[\col!65!black, line width=0.38pt] (\x,7.00) -- (\x,7.20);
}
\node[label, text=cGray!72!black] at (8.78,7.74) {same polymer backbone;\\heterogeneous sulfur environment};

% ============================== C: hero energy landscape ==============================
\node[panel, anchor=west] at (12.00,11.15) {c};
\node[title, anchor=west] at (12.40,11.15) {localized charge traps};
\node[micro, anchor=west, text=cRed!72!black] at (12.40,10.70) {energy landscape of the polymer};
\draw[axis] (12.52,7.06) -- (12.52,10.38);
\node[label, rotate=90] at (12.18,8.72) {energy, $E$};
\draw[cGray!45!black, line width=0.42pt] (12.72,10.05) .. controls (13.42,9.58) and (13.82,9.63) .. (14.26,9.97)
  .. controls (14.74,10.34) and (15.36,10.32) .. (15.82,9.92)
  .. controls (16.24,9.56) and (16.73,9.63) .. (17.32,10.04);
\draw[cBlue!78!black, line width=0.90pt] (13.08,9.86) .. controls (13.34,8.83) and (13.83,8.83) .. (14.09,9.86);
\draw[cRed!82!black, line width=1.05pt] (15.07,9.91) .. controls (15.34,7.58) and (16.20,7.58) .. (16.48,9.91);
\node[circle, fill=cBlue!70, draw=cBlue!90!black, minimum size=1.5mm, inner sep=0pt] at (13.59,9.02) {};
\node[circle, fill=cRed!72, draw=cRed!90!black, minimum size=1.5mm, inner sep=0pt] at (15.77,7.94) {};
\node[label, text=cBlue!80!black, anchor=north] at (13.59,8.57) {shallow\\trap};
\node[label, text=cRed!82!black, anchor=north] at (15.77,7.48) {deep\\trap};
\draw[axis] (13.00,7.00) -- (17.18,7.00);
\node[micro, anchor=north] at (17.18,6.91) {configuration coordinate};
\node[label, text=cGray!70!black] at (14.88,6.36) {trap depth controls\\release time qualitatively};

% ============================== D: transient evidence ==============================
\node[panel, anchor=west] at (0.12,5.56) {d};
\node[title, anchor=west] at (0.52,5.56) {kinetic evidence};
\draw[outline, fill=cGray!7] (0.70,3.66) rectangle (2.12,4.62);
\draw[outline, fill=cAmber!18] (1.04,3.91) rectangle (1.78,4.20);
\draw[signal] (0.42,4.12) -- (1.02,4.12);
\draw[signal] (1.80,4.12) -- (2.42,4.12);
\node[micro, anchor=north] at (1.41,3.56) {MIM test cell};
\draw[axis] (3.00,1.24) -- (5.25,1.24);
\draw[axis] (3.00,1.24) -- (3.00,4.54);
\node[micro, anchor=north] at (5.24,1.15) {$t$};
\node[micro, rotate=90] at (2.62,3.02) {$I(t)$};
\draw[cBlue!82!black, line width=1.0pt] plot[smooth, domain=3.18:5.04, samples=40] (\x,{1.45+2.66*(1+2.1*(\x-3.18))^(-0.72)});
\node[label, fill=white, inner sep=1pt, text=cBlue!82!black] at (4.24,2.37) {$I(t)\propto t^{-n}$};
\node[micro] at (4.18,0.57) {constant bias: current decays};

% ============================== E: ISPD evidence ==============================
\node[panel, anchor=west] at (5.98,5.56) {e};
\node[title, anchor=west] at (6.38,5.56) {surface-potential evidence};
\draw[outline, fill=cAmber!18] (6.45,3.42) rectangle (8.36,3.72);
\draw[outline, fill=cGray!12] (7.23,4.40) rectangle (7.57,4.67);
\draw[outline] (7.40,4.40) -- (7.40,3.78);
\draw[cRed!75!black, line width=0.55pt] (7.40,4.93) -- (7.40,4.68);
\draw[cRed!75!black, line width=0.55pt] (7.25,4.88) -- (7.55,4.88);
\node[micro, anchor=south] at (7.40,5.00) {noncontact probe};
\node[micro, anchor=north] at (7.40,3.30) {charged polymer surface};
\draw[axis] (8.85,1.24) -- (11.18,1.24);
\draw[axis] (8.85,1.24) -- (8.85,4.54);
\node[micro, anchor=north] at (11.18,1.15) {trap energy};
\node[micro, rotate=90] at (8.48,3.02) {$g(E_t)$};
\draw[cTeal!82!black, line width=0.95pt] plot[smooth, domain=9.08:10.92, samples=40] (\x,{1.35+2.35*exp(-7*(\x-10.02)^2)});
\draw[cRed!78!black, line width=0.85pt, dashed] plot[smooth, domain=9.08:10.92, samples=40] (\x,{1.35+1.58*exp(-4*(\x-10.57)^2)});
\node[micro, text=cTeal!85!black] at (9.64,4.18) {shallower};
\node[micro, text=cRed!85!black] at (10.64,3.30) {deeper};

% ============================== F: mechanical evidence ==============================
\node[panel, anchor=west] at (12.00,5.56) {f};
\node[title, anchor=west] at (12.40,5.56) {Coulomb response};
% figure-agent:start panel_f.mechanism_scene
\PanelFFloatingCantilever{redrawF}{15.20,4.55}
\node[micro, anchor=east] at (13.90,1.90) {floating cantilever};
\node[micro, anchor=north, rotate=90] at (16.92,3.05) {electrode};
\node[micro, anchor=south, text=cBlue!78!black] at (16.20,5.18) {$V_{\mathrm{app}}$};
\draw[force] (14.20,2.62) -- (13.18,2.87);
\node[micro, text=cRed!82!black, anchor=south west] at (12.48,2.92) {Coulomb repulsion};
\draw[<->, cGray!65!black, line width=0.38pt] (14.02,1.90) -- (16.50,1.90);
\node[micro, fill=white, inner sep=0.8pt] at (15.26,1.90) {air gap};
% figure-agent:end panel_f.mechanism_scene
\node[label, text=cGray!70!black] at (14.82,0.78) {one grounded voltage-source return;\\sample and cantilever remain floating};

\end{tikzpicture}
\end{document}
